\documentclass[9ps]{beamer}
\usetheme{Frankfurt}
\usecolortheme{seahorse}
\usepackage{amsmath, amsfonts, amssymb, amsthm, graphicx, turnstile, proof,
  logicproof, hyperref, filecontents}

\title{Mechanizing Semi-Axiomatic Sequent Calculi}
\author{Arman Ahmed, Max Gross, Fr\'{e}d\'{e}ric Lahaie-Bertrand}
\date{February 2026}

\begin{document}

\maketitle

\begin{frame}{Goals}
\begin{itemize}
\item Mechanize the Semi-Axiomatic Sequent Calculus (SAX)
  \begin{itemize}
  \item \dots
  \item Proof theory: Admissibility of cut, Cut-elimination, Subformula property
  \item Type theory: Type preservation, Progress, Confluence, Termination 
  \end{itemize}
\item Stretch: Mechanize SNAX 
  \begin{itemize}
  \item ...
  \item Type theory: Type preservation, Progress
  \end{itemize}
\end{itemize}  
\end{frame}

\begin{frame}{SAX: Proof Theory}
\scriptsize
SAX arises from G3 by replacing all \textbf{noninvertible} (e.g. ${\supset}L$)
rules with axioms while leaving all the invertible rules (e.g. ${\supset}R$)
unchanged:
\[\tiny
  \infer[{\supset}L]{\Gamma, A \supset B \vdash C}{
    \Gamma, A \supset B \vdash A \quad
    \Gamma, A \supset B, B \vdash C
  }
  \;\;\Longrightarrow\;\;
  \infer[{\supset}L^0]{\Gamma, A, A \supset B \vdash B}{}
\]
% Example: No cut-free proof of transitivity
\textbf{Cut elimination fails} because axioms introduce formulas that can't be
decomposed:
\[\tiny
  \infer[{\supset}R]{A \supset B, B \supset C \vdash A \supset C}{
    \infer[\text{Cut}_B]{A, A \supset B, B \supset C \vdash C}{
      \infer[{\supset}L^0]{A, A \supset B \vdash B}{} &
      \infer[{\supset}L^0]{B, B \supset C \vdash C}{}
    }
  }
\]
But this cut is \textbf{analytic} $\implies$ subformula property is not lost.

\medskip
\textbf{Snips} are cuts where the cut formula is \emph{eligible}, i.e. it arose
as a subformula inside an axiom's principal formula. Eligibility is tracked by
labeling formula occurrences as eligible ($^*$) or irrelevant ($^0$) and
propagating upward through the derivation:
\[\tiny
  \infer[{\supset}R]{x:A \supset B, y:B \supset C \vdash z:A \supset C}{
    \infer[\text{Snip}_B]{u^*:A, x:A \supset B, y:B \supset C \vdash w^*:C}{
      \infer[{\supset}L^0]{u^*:A, x:A \supset B, y^0:B \supset C \vdash t^*:B}{} &
      \infer[{\supset}L^0]{u^0:A, x^0:A \supset B, t^*:B, y:B \supset C \vdash w^*:C}{}
    }
  }
\]
\textbf{Cut-free in SAX} = no general cuts, snips permitted $\implies$ 
Admissibility of cut
\bigskip
\end{frame}

\begin{frame}{SAX: Shared Memory Interpretation}
    \begin{itemize}
        \item We take our proof terms from the shared memory interpretation as in DeYoung et al.\cite{DeYoung2020}, rather than the message passing interpretation explored by many others.
        \item For a sequent $x_1: a_1, x_2: a_2,\dots, x_3: a_3 \vdash z:C$, we read that the process $P$ (that this sequent represents) reads from memory address $a_1, \dots, a_n$ and writes to address $z$. We can, of course, find $P$ from the deriviation of this sequent. Conversely, the sequent acts as a typing expression for $P$, which we represent $x_1: a_1, x_2: a_2,\dots, x_3: a_3 \vdash P :: (z : C)$.
    \end{itemize}
\end{frame}

\begin{frame}{SAX: Shared Memory Interpretation}
    \begin{itemize}
        \item We find most of the proof terms relatively intuitive, though implication seems to give us some trouble:
    \end{itemize}
    $$\infer[{\supset}L^0]{\Gamma, y^* : A, x: A \supset B \vdash x^R. \langle y, z \rangle :: (z^* : B)}{}$$
    \begin{itemize}
        \item $x$ represents the address of $A \supset B$, $y$ is the standard argument (that the function takes an input of type $A$), $z$ is the address of the result of the function
    \end{itemize}
\end{frame}


\begin{frame}{SAX: Shared Memory Interpretation}
    $$\infer[]{\Gamma \vdash \textsf{case }x^W (\langle y, x \rangle \Rightarrow P) :: (x: A \supset B)}{\Gamma, y: A \vdash P :: (z: B)}$$
    \begin{itemize}
        \item This writes the continuation $\langle y, z \rangle \Rightarrow P$ to the address $x$
    \end{itemize}
\end{frame}

\begin{frame}{Mechanizing SAX}

  Concrete thoughts so far:
    \begin{itemize}
      \item $G_3$ and Theorems 1 and 2 (cut elim. $G_3$) $\checkmark$ (Pfenning 2000) \textcolor{red}{In LF, not Beluga}
      \item SAX $\checkmark$ , Theorem 3 (translation) should follow
      \item Eligibility for snips as well-formedness predicate (a la linearity predicate)
      \item Choosing one of Theorem 4 or 6 (cut elim. in SAX)? Theorem 5 (admiss.) proof looks complicated
      \item Metatheoretic properties look standard
    \end{itemize}
\end{frame}

\begin{frame}
  \scriptsize
\ttfamily

LF hypS : o $\to$ type = ;\\
LF concS : o $\to$ type =\\[4pt]

\quad | initialS : (hypS A $\to$ concS A)\\
\quad | \&RS : concS A $\to$ concS B $\to$ concS (A \& B)\\
\quad | \&L01S : hypS (A \& B) $\to$ concS A\\
\quad | \&L02S : hypS (A \& B) $\to$ concS B\\
\quad | $\supset$RS : (hypS A $\to$ concS B) $\to$ concS (A $\supset$ B)\\
\quad | $\supset$L0S : hypS A $\to$ hypS (A $\supset$ B) $\to$ concS B\\
\quad | $\vee$R01S : hypS A $\to$ concS (A $\vee$ B)\\
\quad | $\vee$R02S : hypS B $\to$ concS (A $\vee$ B)\\
\quad | $\vee$LS : (hypS A $\to$ concS C) $\to$ (hypS B $\to$ concS C) $\to$ (hypS (A $\vee$ B) $\to$ concS C)\\
\quad | cutS : concS A $\to$ (hypS A $\to$ concS C) $\to$ concS C\\
\quad | $\otimes$R0S : hypS A $\to$ hypS B $\to$ concS (A $\otimes$ B)\\
\quad | $\otimes$LS : (hypS A $\to$ hypS B $\to$ concS C) $\to$ (hypS (A $\otimes$ B) $\to$ concS C)\\
\quad | unitRS : concS unit\\
\quad | unitLS : concS C $\to$ (hypS unit $\to$ concS C)\\
;
\end{frame}
\begin{frame} {Related Work}
    \begin{itemize}
        \item Francalaza et. al\cite{FrancalanzaTP24} Presents an implementation of the message-passing interpretation of SAX in Go. Instead of focusing on a pure implementation, we focus on a mechanization of the meta-theoretic properties.
        \item Zhang et al.\cite{zhang2025} gives a mechanization of logical relations for message passing protocols in Coq.
        \item DeYoung et al.\cite{DeYoungP23} presents several metatheoretic proofs for SNAX, including type-safety.
    \end{itemize}

\end{frame}

\begin{frame}{Timeline}
    
\end{frame}

\nocite{*}
\bibliographystyle{plain}
\bibliography{main}

\end{document}